% ID: SEL-D0018-BA-E01
% ESTADO: validada. S019. Revisión matemática, didáctica y visual.
\textbf{a) Determinante.}\par\nopagebreak[4]
Por la primera fila,
\[
\det A=2(5-3)+m(3-m)=4+3m-m^2=(4-m)(m+1).
\]
La inversa existe si $\boxed{m\notin\{-1,4\}}$.

\textbf{b) Resolución para $m=0$.}\par\nopagebreak[4]
Con $m=0$ obtenemos
\[
A=\begin{pmatrix}2&0&0\\1&1&1\\0&3&5\end{pmatrix},
\qquad AA^t=\begin{pmatrix}4&2&0\\2&3&8\\0&8&34\end{pmatrix}.
\]
En $XA=AA^t$ la inversa multiplica por la derecha. También podemos resolver por filas. Para una fila $(u,v,w)$ de $X$, su producto con $A$ es $(2u+v,v+3w,v+5w)$. Si la fila correspondiente de $AA^t$ es $(r,s,t)$, resulta
\[
w=\frac{t-s}{2},\qquad v=s-3w,\qquad u=\frac{r-v}{2}.
\]
Aplicando estas fórmulas a las tres filas,
\[
\boxed{X=\begin{pmatrix}-1/2&5&-1\\13/4&-9/2&5/2\\31/2&-31&13\end{pmatrix}.}
\]
La sustitución en $(2u+v,v+3w,v+5w)$ recupera las tres filas de $AA^t$. La solución es única porque $A$ es invertible.
